\documentclass{letter}
\usepackage{hyperref}

\begin{document}

\begin{letter}{
To the Editors of Mathematical Structures in Computer Science
}

\opening{Dear Editors,}

We are pleased to submit our manuscript titled ``Theory of Incidence -- A Fourth Foundation Beyond Set, Category, and Type'' for consideration for publication in Mathematical Structures in Computer Science.

This work proposes Incidence Theory as a fourth mathematical foundation that unifies Set, Category, and Type theories through a relational primitive. The theory is fully formalized in Lean 4, with machine-checkable proofs establishing its consistency and providing embeddings into existing foundational frameworks.

The key contributions of this work include:
\begin{itemize}
\item A complete axiomatic system (17 axioms) for Incidence Theory with full formalization in Lean 4
\item Relative consistency proof via a ZF set-theoretic model
\item Embeddings of Set, Category, and Type theories into Incidence Theory
\item Linear-algebraic layer connecting logical inference to spectral computation
\item Concrete applications in graph theory, process calculus, and metabolic networks
\end{itemize}

The paper includes extensive formal verification artifacts, making it particularly suitable for MSCS's focus on mathematical structures in computer science. All major theorems are proved in the accompanying Lean 4 codebase, providing unprecedented rigor for a foundational mathematics paper.

The manuscript is approximately 15 pages long, including appendices with Lean code references. It has not been previously published and is not under consideration elsewhere.

We have uploaded the complete Lean 4 formalization to Zenodo (DOI: \href{https://doi.org/10.5281/zenodo.17345516}{10.5281/zenodo.17345516}) and the source code is available at \href{https://github.com/junkawasaki/theory-of-incidence}{https://github.com/junkawasaki/theory-of-incidence}.

We believe this work aligns well with MSCS's mission to publish research on mathematical structures underlying computer science, particularly in the areas of type theory, category theory, and formal methods.

We look forward to the possibility of publishing this work in MSCS.

\closing{Sincerely,}

\end{letter}

\end{document}
